Puisqu'on choisit un élève au hasard parmi les sportifs de haut niveau, on est en situation d'équiprobabilité, et donc les proportions sont assimilables à des probabilités.

	\begin{enumerate}
		\item Voici le tableau complété :

		\begin{tblr}{colspec={l X[c,3cm] X[c,3.3cm] X[c,3cm] X[c,2.5cm]},
			hline{1}={2-Z}{solid},hline{2-Z}={solid},
			vline{1}={2-Z}{solid},vline{2-Z}={solid}}
 			& Seconde & Première & Terminale & Total \\

			Section judo & \textcolor{teal}{$50-40=10$} & \textcolor{teal}{$24-(10+8)=6$} & 8 & 24 \\

			Section aquatique & 40 & \textcolor{teal}{$56-6=50$} & \textcolor{teal}{$14-8=6$} & \textcolor{teal}{$40+50+6=96$} \\

			Total & 50 & \textcolor{teal}{$120-(50+14)=56$} & 14 & 120
		\end{tblr}

		\item $A \cap S$ : "L'élève est en seconde ET en section aquatique";

		$P \big( A \cap S \big) = \dfrac{40}{120}=\dfrac{1}{3}$.

		\item On veut calculer la probabilité $P_S(A)$. Il y a 50 élèves en seconde, donc 40 sont en section aquatique.

		On a donc :\quad $ P_S(A) = \dfrac{40}{50} = \dfrac{4}{5}$.

		\item \begin{enumerate}
			\item Il y a 24 élèves en section Judo :\quad $P \big( J \big)=\dfrac{24}{120}=\dfrac{1}{5}$.

			\item Parmi les 14 élèves de terminale, 8 sont en section judo.\quad $P_T \big( J \big) = \dfrac{8}{14} = \dfrac{4}{7}$.

			\item On remarque que :\quad $P \big( J \big) \neq P_T \big( J \big)$, \quad donc les évènements $J$ et $T$ ne sont pas indépendants.
		\end{enumerate}
	\end{enumerate}
